\documentclass[11pt,aspectratio=169]{beamer}
\usetheme{Boadilla}
\usefonttheme{professionalfonts}
\setbeamertemplate{navigation symbols}{}

\usepackage[T1]{fontenc}
\usepackage{lmodern}
\usepackage{amsmath,amssymb,mathtools}
\usepackage{booktabs}
\usepackage{tabularx}
\usepackage{array}
\usepackage{graphicx}
\usepackage{xcolor}
\usepackage[authoryear,round]{natbib}

\graphicspath{{figs/}}

\newcommand{\good}[1]{\textcolor{blue!70!black}{#1}}
\newcommand{\muted}[1]{\textcolor{gray!75!black}{#1}}

\setbeamercolor{metricbox}{bg=black!5,fg=black}
\setbeamercolor{idbox}{bg=blue!5,fg=black}

\title[Video RDiT]{Video Measurement and Opening-Time RDiT}
\author{}
\date{}

\begin{document}

\begin{frame}{Traffic Outcomes Before and After Opening}
\vspace{-0.5em}
\begin{columns}[T,onlytextwidth]
\column{0.49\textwidth}
\centering
\includegraphics[width=\linewidth,height=0.35\textheight,keepaspectratio,trim=140 50 140 90,clip]{caitian_0709_pre_tracking.jpg}\\[-0.1em]
{\small \textbf{Pre}}
\column{0.49\textwidth}
\centering
\includegraphics[width=\linewidth,height=0.35\textheight,keepaspectratio,trim=140 50 140 90,clip]{caitian_0709_late_post_tracking.jpg}\\[-0.1em]
{\small \textbf{Post}}
\end{columns}

\vspace{0.15em}
\centering
{\scriptsize
\setlength{\tabcolsep}{0.85em}
\renewcommand{\arraystretch}{1.04}
\begin{tabular}{@{}lrrr@{\hspace{1.5em}}rrr@{}}
\toprule
& \multicolumn{3}{c}{\textbf{Flow (vehicles/hour)}}
& \multicolumn{3}{c}{\textbf{Travel time (seconds)}} \\
\cmidrule(lr){2-4}\cmidrule(lr){5-7}
\textbf{Direction} & \textbf{Pre} & \textbf{Post} & \textbf{Change}
& \textbf{Pre} & \textbf{Post} & \textbf{Change} \\
\midrule
$S\to N$ & 2,902 & 3,284 & \good{$+13.2\%$}
& 18.73 & 8.87 & \good{$-52.6\%$} \\
$N\to S$ & 1,484 & 1,712 & \good{$+15.4\%$}
& 3.49 & 3.90 & \good{$+11.8\%$} \\
\bottomrule
\end{tabular}}

\vspace{0.3em}
\begin{beamercolorbox}[wd=0.96\linewidth,sep=0.3em,center]{idbox}
{\scriptsize \textbf{Takeaway:} Flow rises in both directions. Travel time falls only in the capacity-gaining direction.}
\end{beamercolorbox}

\vfill
{\fontsize{4.7}{5.4}\selectfont \muted{\textbf{Data notes.} Pre/Post windows: 0:00--2:45 and 33:30--48:30; deployment is excluded. Flow is normalized to vehicles/hour; travel time is the median matched gate time. Matched tracks (Pre/Post): $S\to N$ 13/655; $N\to S$ 58/403. Source: 0709 outbound video. One-day descriptive comparison, not an RDiT estimate.}}
\end{frame}

\begin{frame}{Identification: Opening-Time RDiT}
\small

\textbf{Design.}
At the scheduled cutoff $c_d$, $S\to N$ capacity shifts from three to four lanes. Actual operation $D_{dt}=Open_{dt}$ is coded from video, and $Y_{dt}\in\{Q_{dt},T_{dt}\}$ measures directional flow or travel time.

\vspace{0.4em}
{\centering
$\displaystyle r_{dt}=t-c_d,\qquad
 Z_{dt}=\mathbf{1}\{r_{dt}\geq0\},\qquad
 D_{dt}=Open_{dt}. $\par}

\vspace{0.5em}
\begin{beamercolorbox}[wd=0.94\linewidth,sep=0.55em,center]{idbox}
{\large
$\displaystyle
Y_{dt}=\alpha_d+\tau_Y D_{dt}
+\beta_-r_{dt}+\beta_+D_{dt}r_{dt}+\varepsilon_{dt},
\qquad |r_{dt}|\leq h.$}
\end{beamercolorbox}

\vspace{0.4em}
\textbf{Estimand.}
$\tau_Y$ is the local effect of opening the extra lane at $c_d$. If operation is delayed, $Z_{dt}$ instruments $D_{dt}$ and $\tau_Y^{FRD}=\pi_Y/\pi_D$.

\vspace{0.35em}
\textbf{Identification.}
\begin{itemize}\small\setlength\itemsep{0.08em}
  \item The scheduled cutoff is predetermined relative to minute-level traffic shocks.
  \item Without opening, potential outcomes evolve smoothly through $c_d$, with no coincident incident or policy change.
\end{itemize}

\vspace{0.05em}
\textbf{Implementation.}
Estimate local linear specifications on repeated dates with one-minute bins, date fixed effects, and date-clustered inference. Exclude deployment; vary $h$; test placebo cutoffs; validate clocks and detection with the reverse camera.

\vfill
{\tiny \muted{Current 0709 estimates are descriptive. Formal RDiT requires repeated opening dates with continuous coverage around the cutoff.}}
\end{frame}

\end{document}
